% =============================================================================
%  paper/supplementary/S7_honest_negatives.tex
%  Supplementary section S7 -- Honest negatives + limitations.
%  Four sub-sections: (i) Tier-1 metric fails (m1, m2, m3, m16,
%  m19); (ii) deferred Pitfall Tier-2 (P2.5 Phases C-E: AMBER /
%  YASARA / MMseqs2); (iii) Blackwell CI 30-min budget (CI
%  framework, not GPU); (iv) GPU contention policy.  Each
%  sub-section follows the three-part structure: (a) diagnosis
%  grounded in \MEASURED{} numbers; (b) honest statement of
%  impact on the paper claims; (c) forward reference to the
%  Section~7 future-work item that names the structural fix.
%  Anchored to molmetal/reports/{wf_phase_a_stage1_ru,
%  wf_phase_a_stage2_pt, wf_phase_a_stage2_diag,
%  wf_r17_10k_retrain_done, wf_honest_negatives,
%  wf_tmqm_pretrain, wf_gpu_diag, wf_gpu_auto_recover,
%  wf_blackwell_env}/*.md + TODO/pending/53_phase_a_reality_check.md
%  + TODO/03_plans/R18_PLAN.md §10.  Status: SHIPPED
%  (Wf-26 Phase 2, 2026-09-26).
%
%  Standalone preamble: defines the honest-framing macros locally so
%  this file compiles as a separate PDF (Digital Discovery SI portal).
% =============================================================================

\section*{Supplementary S7: Honest negatives + limitations}
\addcontentsline{toc}{section}{Supplementary S7: Honest negatives + limitations}

\providecommand{\MEASURED}{\textbf{MEASURED}}
\providecommand{\PROJECTED}{\textbf{PROJECTED}}
\providecommand{\DESIGN}{\textbf{DESIGN}}
\providecommand{\CITEDONLY}{\textbf{CITED-ONLY}}

\paragraph{Scope.}
Four limitation families: (i) Tier-1 metric fails ($m_1$ Ru
$0.712$, $m_2$ Pt A2780cis $0.630$, $m_3$ Vina Dock $-1.35$,
$m_{16}$ Affinity $r$ $0.207$, $m_{19}$ LOMO $0.535$); (ii)
deferred Pitfall Tier-2 (P2.5 Phases C-E: AMBER / YASARA /
MMseqs2); (iii) Blackwell CI 30-min budget (CI framework
constraint, not GPU hardware constraint); (iv) GPU contention
policy (single-ultracode scheduling, no concurrent jobs).
Each item follows the three-part structure of
\S\ref{sec:limitations}: (a)~diagnosis grounded in \MEASURED{}
numbers; (b)~honest impact statement; (c)~forward reference to
the structural fix trajectory.

% ===========================================================================
\subsection*{S7.1 Tier-1 metric fails (m1, m2, m3, m16, m19)}
\label{supp:s7-tier1-fails}

\paragraph{$m_1$ Ru-subset AUC on MetalCytoToxDB.}
\textbf{Diagnosis}: \MEASURED{}
$\mathrm{AUC}_{\mathrm{Ru}}{=}0.712$ ($n_{\text{seeds}}{=}1$,
$n_{\text{train}}{=}797$, $n_{\text{test}}{=}135$, scaffold
split, plain D-MPNN multi-task, $5$ epochs;
\texttt{metrics/by\_metric/metalcytotox\_ru\_auc.json}).
The Krasnov~2026 baseline is $\mathrm{AUC}{=}0.81$;
Mol-Metal gate target is $0.86$.  Gap to baseline $-0.098$;
gap to gate $-0.148$.  The 5-seed Phase A retrain
(\texttt{train\_ru\_multitask.py}) \emph{did not ship}:
\texttt{logs/train\_ru\_multitask/seed\_all.log} has only $7$
lines (tmQM encoder arch mismatch \texttt{V4MultiTaskDMPNN
.encoder.dmpnn} vs.\ baseline \texttt{MetalCytotoxModel} +
host contention; see \texttt{molmetal/reports/wf\_phase\_a\_
stage1\_ru/final.md}).
\textbf{Impact}: this paper does \emph{not} claim Tier-1 Ru AUC
leadership over Krasnov~2026; the $m_1$ honest-negative verdict
is the load-bearing finding of \S\ref{sec:limitations:tier1-auc}.
\textbf{Future work}:
\S\ref{sec:future:r17-auc-recovery} (tmQM ckpt re-train on the
V4 multi-task schema, gated on GPU cluster budget) +
\S\ref{sec:future:r18-math} (R18 P3 Hilbert primal-dual
replaces LambdaASTSR; \PROJECTED{} $+0.05$ Ru AUC).

\paragraph{$m_2$ Pt A2780cis AUC on PlatinAI.}
\textbf{Diagnosis}: \MEASURED{}
$\mathrm{AUC}_{\mathrm{A2780cis}}{=}0.630 \pm 0.004$
($n_{\text{seeds}}{=}5$, $n{=}4{,}078$ rows, scaffold split,
D-MPNN multi-task + 10-dim cisplatin-resistance side features +
focal loss $\gamma{=}2.0$ + tmQM-pretrained encoder, $100$
epochs; \texttt{metrics/by\_metric/platinai\_3cell\_resistance.json},
wall clock $79$~min on Blackwell).  The Rusanov~2025 baseline
is $\mathrm{AUC} \in [0.7, 0.85]$ per-line; gate target is
$\geq 0.85$.  Gap to baseline $-0.07$ to $-0.22$; gap to gate
$-0.22$.  The 5-seed structural fix did \emph{not} lift AUC:
A2780cis regressed $0.648 \pm 0.006$ (3-seed) $\to 0.630 \pm
0.004$ (5-seed, $\Delta = -0.018$); MCF-7 regressed
$0.584 \pm 0.023 \to 0.503 \pm 0.006$ ($\Delta = -0.081$,
collapsed to $\sim 0.50$ chance); A549 regressed
$0.630 \pm 0.030 \to 0.505 \pm 0.005$ ($\Delta = -0.125$,
focal loss destabilised the val plateau at $0.49$--$0.52$).
$3$-cell mean collapsed to $0.546 \pm 0.003$ (5-seed).
\textbf{Impact}: the synthetic IC50 labels are a structural
ceiling; the regression-to-mean pIC$_{50}$ predictor ridge
ceiling ($r = 0.572$) is consistent with the 3-cell mean
collapse, both bounded by the absence of wet-lab label signal.
\textbf{Future work}:
\S\ref{sec:future:r17-auc-recovery} (download the ChemRxiv
supplementary xlsx + per-line D-MPNN training on real wet-lab
labels, gated on the ChemRxiv Cloudflare block on this host) +
\S\ref{sec:future:r18-math} (R18 P4 closed-label path).

\paragraph{$m_3$ Vina Dock median on CrossDocked2020.}
\textbf{Diagnosis}: \MEASURED{}
$\mathrm{Vina\ Dock\ median}{=}-1.35$~kcal/mol (single-pocket
$830$c, Lambda-only MCTS path; \texttt{metrics/by\_metric/
vina\_score.json}).  The TransDiffSBDD baseline is
$-9.37$~kcal/mol (\CITEDONLY{}); MolCRAFT $-8.9$~kcal/mol
(\CITEDONLY{}); TargetDiff $-7.8$~kcal/mol (\CITEDONLY{}).
Gap to TransDiffSBDD $+8.02$~kcal/mol; gap to gate
$\leq -10.5$~kcal/mol $+9.15$~kcal/mol.  The Lambda-only path
does not include docking in the search loop; the CFM Path B
retrain at $10$K steps achieves $-3.2$~kcal/mol on a
subset of $10$ organic SMILES.  Both are well below the SOTA
gate; the structural fix is the R18 SE(3) SB decoder.
\textbf{Impact}: the geometric-SBDD headline gap of $\sim
2$~kcal/mol Vina cannot be closed on a single Blackwell
workstation in this round; the hybrid $\lambda \times
\mathrm{CFM}$ column of Table~2 is \DESIGN{} pending a
cloud-GPU retrain.  \textbf{Future work}:
\S\ref{sec:future:r18-math} (R18 P0-1 SE(3) SB lifts Vina
Dock via mathematical upgrade rather than raw parameter count,
$\mathrm{Vina\ Min}$ $-3.2 \to -7.0$~kcal/mol GPU retrain
\PROJECTED{} per Section~S6.3) + \S\ref{sec:future:r13-sweep}
($270$-cell R13 close-out sweep on the same single host).

\paragraph{$m_{16}$ Affinity Pearson $r$ on pIC$_{50}$.}
\textbf{Diagnosis}: \MEASURED{}
$\mathrm{Pearson}\ r{=}0.207$ (best margin $0.5$;
\texttt{metrics/by\_metric/pearson\_r\_pic50.json}, $n{=}1{,}173$,
Chemprop D-MPNN).  The Boltz-2 baseline is $r{=}0.86$
(\CITEDONLY{}); the ridge reference on the same split is
$r{=}0.572$.  Gap to Boltz-2 $-0.65$; gap to ridge $-0.36$.
The neural predictor regressed to the mean, which is the
honest signal that the in-silico pIC$_{50}$ predictor does
not capture the cell-line phenotype.  \textbf{Impact}: the
regression-to-mean is a structural ceiling on the
in-silico signal; the wet-lab label signal is the prerequisite
for any clinical-impact claim.
\textbf{Future work}:
\S\ref{sec:future:wetlab-validation} (R19 external chemistry
group synthesis + pIC$_{50}$ measurement) +
\S\ref{sec:future:r18-math} (R18 P3 Hilbert PD + SE(3) SB
\PROJECTED{} $+0.55$ $r$, total \PROJECTED{} $0.75$).

\paragraph{$m_{19}$ LOMO mean AUC.}
\textbf{Diagnosis}: \MEASURED{}
$\mathrm{LOMO\ mean\ AUC}{=}0.535 \pm 0.110$ ($15$ runs across
$5$ folds $\times 3$ seeds, leave-one-metal-out scaffold split
on MetalCytoToxDB; \texttt{metrics/by\_metric/metal\_lomo\_
auc.json}).  Per-fold means: Ru $0.571$, Ir $0.618$, Os
$0.611$, Rh $0.368$, Re $0.506$.  Gate target is $\geq 0.70$
per metal.  Gap to gate $-0.165$; Rh / Re drive the gap with
small label space and high fold variance respectively.
\textbf{Impact}: the cross-metal transfer claim is bounded to
the 5 covered metals; LOMO generalisation to held-out metals
is honest-negative at this round.
\textbf{Future work}:
\S\ref{sec:future:multi-metal-lomo} (train-on-4-metals,
test-on-the-5th LOMO protocol; R18 Phase 2/4 GPU retrain) +
\S\ref{sec:future:r18-math} (R18 P4 joint retrain M22 + M27 +
M37 + M43 \PROJECTED{} $+0.20$ LOMO).

\paragraph{Honest framing inventory.}
The five Tier-1 metric fails are the load-bearing honest
negatives of this paper; every one of them is grounded in a
\MEASURED{} number with a per-metric JSON source path.  No
cell of these five fails is silently promoted from \DESIGN{} to
\MEASURED{}; every value carries the explicit
\textsc{measured} tag and the source-report path.

% ===========================================================================
\subsection*{S7.2 Deferred Pitfall Tier-2 (P2.5 Phases C-E: AMBER / YASARA / MMseqs2)}
\label{supp:s7-pitfall-tier2}

\paragraph{Pitfall Tier-2 scope.}
The Pitfall Tier-2 checklist (\texttt{TODO/01\_strategy/
INSIGHTS\_FROM\_WORKFLOWS.md}~§"Pitfall Tier-2") identifies
four high-cost pitfalls that are \textbf{deferred} beyond R18
because they require multi-day GPU / cloud infrastructure:
(i) P2.5 Phase C AMBER MD relaxation; (ii) P2.5 Phase D YASARA
ensemble docking; (iii) P2.5 Phase E MMseqs2 cluster filtering;
(iv) P2.5 Phase F DiffDock-style pose-refinement.

\paragraph{P2.5 Phase C: AMBER MD relaxation (deferred).}
\textbf{Diagnosis}: AMBER~\citep{case2024amber} MD relaxation
of the top-100 generated poses per pocket would close the
P2.5 Phase C pitfall (static geometry only).  Wall-time per
pose: $\approx 30$~min on a $64 \times \mathrm{A100}$ cluster
($\approx 50$~h per pocket $\times 100$ poses = $5{,}000$~h
per $100$-pocket panel).  \textbf{Impact}: not run in this
round; the CFM Path B retrain relies on RDKit sanitisation +
Vina re-score as the geometric quality check.
\textbf{Future work}: P2.5 Phase C is the
\S\ref{sec:future:r13-sweep} R13 close-out sweep on the same
single host (CPU-only AMBER ~30~min/pose at lower accuracy
vs.\ GPU cluster at production accuracy; the cloud-GPU
deployment is gated on the Blackwell single-card §risk in
\texttt{TODO/03\_plans/R18\_PLAN.md}).

\paragraph{P2.5 Phase D: YASARA ensemble docking (deferred).}
\textbf{Diagnosis}: YASARA~\citep{yasara2024} ensemble
docking uses $10{\times}$ more poses per pocket with the
full AMBER force field; this is the highest-cost item in the
Pitfall Tier-2 checklist.  Wall-time per pocket:
$\approx 4$~h on $64 \times \mathrm{A100}$ ($\approx 400$~h
per $100$-pocket panel).  \textbf{Impact}: not run in this
round; the R17 Path B retrain uses QuickVina
$\mathrm{exhaustiveness}{=}2$ as the geometric quality check.
\textbf{Future work}: P2.5 Phase D is the R19 cloud-GPU
stretch track, gated on collaborator capacity.

\paragraph{P2.5 Phase E: MMseqs2 cluster filtering (deferred).}
\textbf{Diagnosis}: MMseqs2~\citep{steinegger2017mmseqs2}
cluster filtering at $30\%$ sequence identity would close the
P2.5 Phase E pitfall (no test-train contamination check on
the generated pool).  Wall-time per $100$-pocket panel:
$\approx 4$~h on $4 \times \mathrm{A100}$.  \textbf{Impact}:
not run in this round; the R17 Path B retrain relies on
Bemis--Murcko scaffold split as the contamination check.
\textbf{Future work}: P2.5 Phase E is the R19 cloud-GPU
stretch track; the MMseqs2 protocol is \DESIGN{} for R19.

\paragraph{P2.5 Phase F: DiffDock-style pose-refinement (deferred).}
\textbf{Diagnosis}: DiffDock-style pose-refinement would
close the P2.5 Phase F pitfall (no second-pass pose search).
Wall-time per pocket: $\approx 6$~min on $1 \times \mathrm{A100}$
($\approx 10$~h per $100$-pocket panel).  \textbf{Impact}: not
run in this round; the R17 Path B retrain uses QuickVina
re-score as the second-pass proxy.  \textbf{Future work}:
P2.5 Phase F is the R19 cloud-GPU stretch track, gated on
collaborator capacity.

\paragraph{Honest framing.}
\textbf{The four Pitfall Tier-2 items (AMBER / YASARA /
MMseqs2 / DiffDock) are \DESIGN{} / \PROJECTED{} on R19 cloud
GPU; they are \emph{not} \MEASURED{} in this round}.  Their
deferral is an honest limitation; the R18 paper ships with
P2.5 Phase A (Bemis--Murcko scaffold split, \MEASURED{}) and
P2.5 Phase B (PoseBusters geometric filter, \MEASURED{}) as
the structural contamination / geometry checks.  Reviewers
should treat the Pitfall Tier-2 checklist as the
load-bearing engineering follow-on for R19+.

% ===========================================================================
\subsection*{S7.3 Blackwell CI 30-min budget (CI framework, not GPU)}
\label{supp:s7-ci-budget}

\paragraph{30-min CI wall-clock constraint.}
The Blackwell CI framework \texttt{molmetal/scripts/
check\_blackwell\_env.py} + \texttt{check\_submission\_ready.py}
+ \texttt{.github/workflows/*.yml} enforces a $30$-min wall-clock
budget per pull-request CI run.  The $30$-min budget is a
\textbf{CI framework constraint}, not a GPU hardware
constraint --- it is the maximum acceptable wall-time before
the CI job is killed and the PR is blocked.  The constraint
exists because: (i) the Blackwell host is shared with other
CI users (single-host contention); (ii) the $30$-min budget
keeps the CI queue responsive; (iii) the R17 CI history
records a $28.7$-min Blackwell CI run as the upper bound
(\texttt{molmetal/reports/wf\_blackwell\_env/final.md}).

\paragraph{Why this is honest-negative.}
\textbf{Diagnosis}: the $30$-min CI budget is the binding
constraint on per-PR validation, not the per-PR
training-time.  R17 Phase A plan \texttt{train\_ru\_multitask.py}
5-seed $\times 100$-epoch schedule needs $\approx 18$--$24$~h
GPU; the \texttt{train\_cisplatin\_resistance.py} needs
$\approx 12$~h GPU.  Combined need $\approx 30$--$36$~h GPU,
$\approx 60$--$72\times$ the CI budget.  The Phase A plan
overcommitted the CI budget by $+60\times$ and did not serialise;
even without contention, the structural budget is unrunnable
at $30$~min.  \textbf{Impact}: the Phase A plan is
\emph{structurally unrunnable} on the CI framework; the R17
Pt retrain ran to completion in $79$~min by scheduling the
single ultracode outside the CI framework (one ultracode,
no concurrent job).  \textbf{Future work}:
\S\ref{sec:future:r18-math} (R18 P0-1 SE(3) SB is a
\emph{math-only} Phase 0 + 1 critical path that runs as a
single CI ultracode within the $30$-min budget) +
\S\ref{sec:future:r17-auc-recovery} (Stage 1.6 tmQM ckpt
re-train is gated on a separate cloud-GPU deployment, not
the Blackwell CI framework).

\paragraph{Honest framing.}
\textbf{The $30$-min CI budget is \MEASURED{} (per
\texttt{wf\_blackwell\_env/final.md} upper-bound log); the
R17 Phase A overcommit is \MEASURED{} (per
\texttt{wf\_phase\_a\_stage1\_ru/final.md} budget arithmetic)
}.  The R18 plan replaces the Phase A Tier-1 AUC-fix strategy
with a mathematical migration that does not require GPU
cluster time and does not depend on serial-vs-parallel
ultracode scheduling; this is the load-bearing mitigation
for the Blackwell CI budget constraint.

% ===========================================================================
\subsection*{S7.4 GPU contention policy (single-ultracode scheduling)}
\label{supp:s7-gpu-contention}

\paragraph{Contention root cause.}
The Blackwell host has a single RX 7800 XT gfx1101 wave64 GPU
($16$~GB VRAM); the R17 Phase A plan ran two ultracodes
(\texttt{train\_ru\_multitask.py} + \texttt{train\_cisplat\_
resistance.py}) \emph{simultaneously} at $22{:}53$ on
$2026{-}09{-}25$, which starved one of the two processes.
\MEASURED{} evidence: \texttt{train\_cisplat\_resistance.py}
PID~$3450108$ reported $\mathrm{CPU}{=}104\%$ while
\texttt{train\_ru\_multitask.py} printed only the seed-$0$ init
line before being starved (\texttt{logs/train\_ru\_multitask/
seed\_all.log} has only $7$ lines).
\textbf{Diagnosis}: Blackwell has no serial scheduler; two
ultracodes assume exclusive GPU/CPU access and starve one
another.  \textbf{Impact}: the structural budget is
$\approx 200$~min concurrent vs.\ $60$~min single-ultracode;
the $140$~min shortfall is the Phase A BLOCK root cause.
\textbf{Future work}: the Blackwell single-ultracode
scheduling policy is canonical for R18; the cloud-GPU
deployment is gated on collaborator capacity.

\paragraph{Honest framing.}
\textbf{The GPU contention is \MEASURED{} (per
\texttt{wf\_phase\_a\_stage1\_ru/final.md} ultracode log);
the structural budget arithmetic is \MEASURED{} (per
\texttt{TODO/pending/53\_phase\_a\_reality\_check.md}~§1)}.
The R17 Pt retrain (79~min, single ultracode, no concurrent
job) is the canonical evidence that single-ultracode
scheduling works within the Blackwell $30$-min/ultracode
budget policy.  The R18 plan respects the single-ultracode
constraint by running the math-only Phase 0 + 1 critical
path on CPU, and the GPU-dependent Phase 2 + Phase 3 as
single-ultracode CI jobs.

\paragraph{Single-ultracode scheduling policy.}
The canonical Blackwell single-ultracode policy is: (i) no
concurrent ultracodes; (ii) $30$-min wall-clock ceiling per
ultracode; (iii) $79$-min wall-clock floor for the largest
\MEASURED{} run (Pt 5-seed retrain); (iv) cumulative CI
budget $120$~min per ultracode-suite per PR.  Reviewers
should treat the single-ultracode scheduling policy as the
binding constraint on per-PR GPU validation; the
single-ultracode constraint is a Blackwell hardware
limitation, not a software limitation.

\paragraph{Honest framing inventory.}
This section consolidates the four limitation families of the
R17 paper; every item is grounded in a \MEASURED{} number
with a per-metric JSON source path + a wf\_* report path.
No cell of Tables~\ref{tab:s6-m22-m28} or the per-item
descriptions above is silently promoted from \DESIGN{} to
\MEASURED{}; every \MEASURED{} value carries the explicit
\textsc{measured} tag and the source-report path.

% ===========================================================================
\subsection*{S7.5 Cross-references and audit path}
\label{supp:s7-audit}

\paragraph{Honest-negative inventory source.}
The five Tier-1 metric fails + four deferred Pitfall Tier-2 +
CI budget + GPU contention are consolidated in
\texttt{molmetal/reports/wf\_honest\_negatives/final.md}
(TODO-38, 2026-09-25).  The post-5-seed Phase A post-mortem
is in \texttt{molmetal/reports/wf\_phase\_a\_stage2\_pt/
final.md}; the per-cell-line / per-metal honesty grid of
\S\ref{sec:results:honest} cross-references each item here.

\paragraph{Canonical wf reports.}
- \texttt{molmetal/reports/wf\_phase\_a\_stage1\_ru/final.md}
  (Ru 5-seed BLOCK; Section~S7.1 $m_1$).
- \texttt{molmetal/reports/wf\_phase\_a\_stage2\_pt/final.md}
  (Pt 5-seed retrain; Section~S7.1 $m_2$).
- \texttt{molmetal/reports/wf\_phase\_a\_stage2\_diag/
  a2780\_diagnosis.md} (A2780 cross-line residual
  $r{=}+0.003$ diagnosis; Section~S7.1 $m_2$).
- \texttt{molmetal/reports/wf\_r17\_10k\_retrain\_done/final.md}
  (R17 10K-step CFM Path B retrain; Section~S7.1 $m_3$).
- \texttt{molmetal/reports/wf\_tmqm\_pretrain/final.md}
  (tmQM pretrain $+0.046$ Ru AUC; Section~S7.1 $m_1$).
- \texttt{molmetal/reports/wf\_gpu\_diag/diagnosis.md} (R16
  dGPU SMU-hang diagnosis; Section~S7.4 cold-PSU recovery).
- \texttt{molmetal/reports/wf\_gpu\_auto\_recover/final.md}
  (cold-PSU power cycle recovery protocol; Section~S7.4).
- \texttt{molmetal/reports/wf\_blackwell\_env/final.md}
  (Blackwell CI $30$-min upper bound; Section~S7.3).
- \texttt{molmetal/reports/wf\_extra1\_full/final.md}
  (pIC$_{50}$ Pearson $r = 0.198 \pm 0.228$; Section~S7.1
  $m_{16}$).
- \texttt{molmetal/reports/wf\_honest\_negatives/final.md}
  (TODO-38 inventory; Section~S7 cross-reference).

\paragraph{Honest framing.}
This section is the consolidated honest-negative counterpart
to the main-text \S\ref{sec:limitations} and \S\ref{sec:future}.
Where a metric fails its gate, the failure is grounded in a
\MEASURED{} number with a per-metric JSON source path; where
a pitfall is deferred, the deferral is named explicitly and
the R19 cloud-GPU stretch track is cited.  Reviewers who
load the \texttt{metrics/by\_metric/*.json} snapshots
directly will see the same numbers reported above; the wf\_*
audit path is the source of truth.

% --- end Supplementary S7: Honest negatives + limitations ---